% ============================================================
%  Chapitre 4 — Analyse Factorielle des Correspondances (AFC)
% ============================================================
\section{Analyse Factorielle des Correspondances (AFC)}

% ────────────────────────────────────────────────────────────
\subsection{4.1 Position du problème}
% ────────────────────────────────────────────────────────────
\begin{frame}{4.1 — Pourquoi l'AFC ?}
\begin{columns}[T]
\column{0.55\textwidth}
\textbf{Données} : table de contingence $\mathbf{N} = (n_{ij}) \in \N^{I \times J}$
croisant deux variables \alert{qualitatives}.\\[0.1cm]

\textbf{Objectifs}
\begin{enumerate}
  \item Décrire les associations entre modalités lignes et colonnes
  \item Tester / quantifier l'\textbf{écart à l'indépendance}
  \item Représenter \textbf{simultanément} lignes et colonnes
        sur un plan factoriel
\end{enumerate}

\vspace{0.2cm}
\textbf{Origine} : Benzécri (1973), école française.
\textbf{AFC} = ACP \emph{généralisée} avec :
\begin{itemize}\small
  \item Tableau de profils
  \item Métrique du $\chi^2$
\end{itemize}
\column{0.42\textwidth}
\centering
\footnotesize
\textbf{CSP $\times$ Diplôme} (INSEE)
\renewcommand{\arraystretch}{1.05}
\begin{tabular}{l|cccc|c}
\toprule
 & Bac & Lic. & Mast. & Doct. & total \\
\midrule
Cadre   & 12 & 65 & 180 & 95 & 352 \\
ProfInt & 88 & 110 & 50  & 5  & 253 \\
Employé & 230& 45 & 12  & 1  & 288 \\
Ouvrier & 195& 8  & 2   & 0  & 205 \\
\midrule
total   & 525& 228& 244 & 101& 1098\\
\bottomrule
\end{tabular}

\vspace{0.4cm}
\begin{block}{\small Question}
\small Quels diplômes sont \emph{associés} à quelles CSP ?
\end{block}
\end{columns}
\end{frame}

% ────────────────────────────────────────────────────────────
\subsection{4.2 Profils lignes et colonnes}
% ────────────────────────────────────────────────────────────
\begin{frame}{4.2 — Profils lignes et profils colonnes}
\begin{definition}[Tableau de fréquences relatives]
\[
f_{ij} = \frac{n_{ij}}{n}, \qquad
f_{i \bullet} = \sum_j f_{ij}, \qquad
f_{\bullet j} = \sum_i f_{ij}
\]
où $n = \sum_{i,j} n_{ij}$ est l'effectif total.
\end{definition}

\begin{definition}[Profil ligne, profil colonne]
\begin{itemize}
  \item \textbf{Profil de la ligne $i$} :
        $\mathbf{L}_i = \left(\dfrac{f_{i1}}{f_{i\bullet}}, \ldots, \dfrac{f_{iJ}}{f_{i\bullet}}\right) \in \R^J$
  \item \textbf{Profil de la colonne $j$} :
        $\mathbf{C}_j = \left(\dfrac{f_{1j}}{f_{\bullet j}}, \ldots, \dfrac{f_{Ij}}{f_{\bullet j}}\right) \in \R^I$
\end{itemize}
Chaque profil est une \alert{distribution conditionnelle} (somme $= 1$).
\end{definition}

\begin{rema}[Profil moyen]
\small Le \textbf{profil moyen ligne} est $(f_{\bullet 1}, \ldots, f_{\bullet J})$ (la marge colonne).
Sous l'\textbf{indépendance} $\chi^2$, tous les profils lignes seraient égaux à ce profil moyen.
\end{rema}
\end{frame}

% ────────────────────────────────────────────────────────────
\subsection{4.3 Distance du chi-deux}
% ────────────────────────────────────────────────────────────
\begin{frame}{4.3 — Distance du $\chi^2$}
\begin{definition}[Distance du $\chi^2$ entre profils lignes]
\[
d_{\chi^2}^2(\mathbf{L}_i, \mathbf{L}_{i'})
= \sum_{j=1}^J \frac{1}{f_{\bullet j}}
  \left(\frac{f_{ij}}{f_{i\bullet}} - \frac{f_{i'j}}{f_{i'\bullet}}\right)^2.
\]
C'est une distance euclidienne dans $\R^J$ avec \alert{métrique
$\mathbf{M} = \mathbf{D}_{1/f_{\bullet j}}$} (inverses des fréquences marginales).
\end{definition}

\begin{prop}[Propriétés essentielles]
\begin{enumerate}
  \item \textbf{Pondération} : les colonnes \emph{rares} pèsent plus lourd
        ($1/f_{\bullet j}$).
  \item \textbf{Équivalence distributionnelle} : si deux colonnes sont
        proportionnelles, leur fusion ne change pas les distances entre lignes.
  \item \textbf{Symétrie} : on définit de même la distance entre profils
        colonnes avec $\mathbf{M} = \mathbf{D}_{1/f_{i\bullet}}$.
\end{enumerate}
\end{prop}

\begin{rema}
\small La métrique du $\chi^2$ \alert{équilibre l'influence} des modalités peu
fréquentes mais informatives.
\end{rema}
\end{frame}

% ────────────────────────────────────────────────────────────
\subsection{4.4 Lien avec le test d'indépendance}
% ────────────────────────────────────────────────────────────
\begin{frame}{4.4 — Test d'indépendance et inertie totale}
\textbf{Hypothèse d'indépendance} : $f_{ij} = f_{i\bullet} \, f_{\bullet j}$.

\begin{theorem}[Statistique du $\chi^2$ et inertie]
La statistique de Pearson est :
\[
\chi^2 = n \sum_{i,j} \frac{(f_{ij} - f_{i\bullet} f_{\bullet j})^2}{f_{i\bullet} f_{\bullet j}}
       = n \, \Phi^2.
\]
Le coefficient $\Phi^2 = \chi^2 / n$ est exactement l'\alert{inertie totale} du
nuage des profils lignes (ou colonnes), pondérés par leurs marges :
\[
\boxed{\;
I_{\text{totale}}
= \sum_{i=1}^I f_{i\bullet}\, d_{\chi^2}^2(\mathbf{L}_i, \mathbf{L}_{\bullet})
= \Phi^2 = \chi^2 / n.
\;}
\]
\end{theorem}

\begin{rema}[Conséquence]
\small L'AFC \emph{décompose} la statistique du $\chi^2$ : chaque axe factoriel
explique une part de l'écart à l'indépendance.
\[
\Phi^2 = \sum_{k=1}^{r} \lambda_k, \qquad r = \min(I-1, J-1).
\]
\end{rema}
\end{frame}

% ────────────────────────────────────────────────────────────
\subsection{4.5 Formulation et résolution}
% ────────────────────────────────────────────────────────────
\begin{frame}{4.5 — Formulation comme SVD généralisée}
\textbf{Tableau centré} : $\widetilde{f}_{ij} = f_{ij} - f_{i\bullet} f_{\bullet j}$.

\begin{theorem}[Diagonalisation]
Soit $\mathbf{S}$ la matrice de terme général $\widetilde{f}_{ij}/\sqrt{f_{i\bullet} f_{\bullet j}}$.
La SVD $\mathbf{S} = \mathbf{U}\mathbf{\Sigma}\mathbf{V}^\top$ fournit les axes principaux
de l'AFC :
\begin{itemize}
  \item Valeurs propres : $\lambda_k = \sigma_k^2$ (vérifient $\lambda_k \leq 1$)
  \item Inertie totale : $\Phi^2 = \sum_k \lambda_k$
  \item Axes des lignes : $\mathbf{D}_{1/\sqrt{f_{i\bullet}}} \mathbf{U}$
  \item Axes des colonnes : $\mathbf{D}_{1/\sqrt{f_{\bullet j}}} \mathbf{V}$
\end{itemize}
\end{theorem}

\begin{block}{Cadre unificateur}
\small AFC $=$ ACP du triplet $(\mathbf{X}, \mathbf{M}, \mathbf{D}_p)$ avec :
\begin{itemize}\small
  \item $\mathbf{X}$ = tableau des profils lignes
  \item $\mathbf{M} = \mathbf{D}_{1/f_{\bullet j}}$ (métrique $\chi^2$)
  \item $\mathbf{D}_p = \mathbf{D}_{f_{i\bullet}}$ (pondération par marges)
\end{itemize}
\end{block}
\end{frame}

% ────────────────────────────────────────────────────────────
\subsection{4.6 Coordonnées factorielles}
% ────────────────────────────────────────────────────────────
\begin{frame}{4.6 — Coordonnées factorielles}
\begin{definition}[Coordonnées des lignes et des colonnes]
Sur le $k$-ième axe, les coordonnées sont :
\[
F_k(i) \quad\text{(coordonnée de la ligne } i\text{)}, \qquad
G_k(j) \quad\text{(coordonnée de la colonne } j\text{)}.
\]
\end{definition}

\begin{theorem}[Inertie expliquée par axe]
\[
\sum_{i=1}^I f_{i\bullet}\, F_k(i)^2 = \lambda_k,
\qquad
\sum_{j=1}^J f_{\bullet j}\, G_k(j)^2 = \lambda_k.
\]
Lignes et colonnes contribuent \emph{ensemble} à la même inertie $\lambda_k$.
\end{theorem}

\begin{rema}[Double point de vue]
\small L'AFC traite \alert{symétriquement} lignes et colonnes : on diagonalise
deux matrices semblables qui partagent les mêmes valeurs propres non nulles
(dualité). C'est ce qui justifie la \textbf{représentation simultanée}.
\end{rema}
\end{frame}

% ────────────────────────────────────────────────────────────
\subsection{4.7 Formules de transition}
% ────────────────────────────────────────────────────────────
\begin{frame}{4.7 — Formules barycentriques (transition)}
\begin{theorem}[Relations de transition]
Pour tout axe $k$ et avec $\lambda_k$ la valeur propre associée :
\[
\boxed{\;
F_k(i) = \frac{1}{\sqrt{\lambda_k}}
         \sum_{j=1}^J \frac{f_{ij}}{f_{i\bullet}}\, G_k(j),
\qquad
G_k(j) = \frac{1}{\sqrt{\lambda_k}}
         \sum_{i=1}^I \frac{f_{ij}}{f_{\bullet j}}\, F_k(i).
\;}
\]
\end{theorem}

\begin{block}{Lecture barycentrique}
\small La position d'une ligne $i$ est, à un facteur près, le \alert{barycentre}
des positions des colonnes $j$ pondérées par son profil $f_{ij}/f_{i\bullet}$.
Et inversement.
\end{block}

\begin{rema}[Interprétation graphique]
\small Une modalité ligne et une modalité colonne \alert{proches} sur le plan
factoriel sont dans une \textbf{relation d'attraction} : leur fréquence
conjointe dépasse l'attendu sous indépendance.
\end{rema}
\end{frame}

% ────────────────────────────────────────────────────────────
\subsection{4.8 Représentation simultanée}
% ────────────────────────────────────────────────────────────
\begin{frame}{4.8 — Représentation simultanée}
\begin{columns}[T]
\column{0.49\textwidth}
\textbf{Biplot AFC}\\
\small Sur le \emph{même} plan, on représente :
\begin{itemize}\small
  \item les modalités \alert{lignes} ($I$ points)
  \item les modalités \alert{colonnes} ($J$ points)
\end{itemize}

\textbf{Règles de lecture}
\begin{itemize}\small
  \item Deux \emph{lignes} proches : profils similaires
  \item Deux \emph{colonnes} proches : profils similaires
  \item Ligne $\times$ colonne proche : association forte
  \item Origine $=$ profil moyen
        (proche $=$ proche de l'indépendance)
\end{itemize}
\column{0.49\textwidth}
\begin{tikzpicture}[scale=0.95]
\draw[->, gray!50] (-2.5,0) -- (2.5,0) node[right, font=\tiny]{$F_1$};
\draw[->, gray!50] (0,-2) -- (0,2) node[above, font=\tiny]{$F_2$};

\fill[deepblue]   ( 1.7, 1.5)  circle(2.5pt) node[above, font=\tiny] {Cadre};
\fill[deepblue]   ( 0.4, 0.6)  circle(2.5pt) node[above, font=\tiny] {ProfInt};
\fill[deepblue]   (-1.0,-0.5)  circle(2.5pt) node[below, font=\tiny] {Employé};
\fill[deepblue]   (-2.0,-1.4)  circle(2.5pt) node[below, font=\tiny] {Ouvrier};

\fill[deeporange] (-2.0,-1.0)  circle(2.5pt) node[above, font=\tiny] {Bac};
\fill[deeporange] (-0.3, 0.4)  circle(2.5pt) node[right, font=\tiny] {Lic.};
\fill[deeporange] ( 1.2, 1.2)  circle(2.5pt) node[right, font=\tiny] {Mast.};
\fill[deeporange] ( 1.9, 1.6)  circle(2.5pt) node[right, font=\tiny] {Doct.};

\node[font=\tiny, gray, align=center] at (-2.0,1.7) {Lignes (CSP)\\Colonnes (Diplôme)};
\end{tikzpicture}
\end{columns}
\end{frame}

% ────────────────────────────────────────────────────────────
\subsection{4.9 Aides à l'interprétation}
% ────────────────────────────────────────────────────────────
\begin{frame}{4.9 — Aides à l'interprétation}
\begin{definition}[Qualité de représentation et contributions]
Pour la ligne $i$ sur l'axe $k$ :
\[
\cos^2_k(i) = \frac{F_k(i)^2}{\sum_{l \geq 1} F_l(i)^2},
\qquad
\text{ctr}_k(i) = \frac{f_{i\bullet}\, F_k(i)^2}{\lambda_k}.
\]
Définitions analogues pour les colonnes.
\end{definition}

\begin{block}{Règles de lecture}
\begin{itemize}\small
  \item \alert{Contribution élevée} : la modalité \emph{construit} l'axe
        (seuil pratique : ctr $> 1/I$ ou $1/J$)
  \item \alert{$\cos^2$ élevé} : la modalité est \emph{bien représentée}
  \item Une modalité avec $\cos^2$ faible doit être interprétée avec prudence
\end{itemize}
\end{block}

\begin{rema}
\small On peut décomposer le $\chi^2$ en sommes par cellule (\textbf{contribution
au $\chi^2$}) :
$\chi^2_{ij} = n \cdot (f_{ij} - f_{i\bullet}f_{\bullet j})^2 / (f_{i\bullet}f_{\bullet j})$,
ce qui complète l'AFC.
\end{rema}
\end{frame}

% ────────────────────────────────────────────────────────────
\subsection{4.10 Reconstitution du tableau}
% ────────────────────────────────────────────────────────────
\begin{frame}{4.10 — Reconstitution et approximation}
\begin{theorem}[Formule de reconstitution]
Le tableau de fréquences peut être reconstruit à partir des $r$ axes
principaux :
\[
f_{ij} = f_{i\bullet} f_{\bullet j}
        \left(1 + \sum_{k=1}^r \frac{F_k(i)\, G_k(j)}{\sqrt{\lambda_k}}\right).
\]
\end{theorem}

\begin{block}{Lecture}
\begin{itemize}\small
  \item Le terme $f_{i\bullet} f_{\bullet j}$ est l'\alert{indépendance théorique}.
  \item La somme corrective décompose les écarts à l'indépendance par axe.
  \item Tronquer la somme à $q < r$ axes donne la \textbf{meilleure
        approximation de rang $q$} (Eckart-Young).
\end{itemize}
\end{block}

\begin{rema}[Comparaison à l'ACP]
\small AFC et ACP partagent la même structure : SVD généralisée + Eckart-Young.
La différence est uniquement dans le \emph{triplet} $(\mathbf{X}, \mathbf{M}, \mathbf{D}_p)$.
\end{rema}
\end{frame}

% ────────────────────────────────────────────────────────────
\subsection{4.11 Modalités supplémentaires}
% ────────────────────────────────────────────────────────────
\begin{frame}{4.11 — Modalités et tableaux supplémentaires}
\textbf{Modalité supplémentaire} : projetée sur les axes \emph{après} construction,
sans contribuer à leur calcul.

\begin{block}{Formules de projection}
Pour une ligne supplémentaire $i^*$ de profil $\mathbf{L}_{i^*}$ :
\[
F_k(i^*) = \frac{1}{\sqrt{\lambda_k}}
         \sum_{j=1}^J \frac{f_{i^*j}}{f_{i^*\bullet}}\, G_k(j).
\]
Symétriquement pour une colonne supplémentaire.
\end{block}

\begin{columns}[T]
\column{0.49\textwidth}
\textbf{Cas d'usage}
\begin{itemize}\small
  \item Comparer un échantillon test à une population de référence
  \item Suivre une modalité dans le temps (cohorte)
  \item Caractériser un axe par une variable externe
\end{itemize}
\column{0.49\textwidth}
\begin{alertblock}{Attention}
\small Une modalité supplémentaire \emph{rare} ($f \to 0$) peut être très
mal positionnée : on l'interprète avec son $\cos^2$.
\end{alertblock}
\end{columns}
\end{frame}

% ────────────────────────────────────────────────────────────
\subsection{4.12 Code R}
% ────────────────────────────────────────────────────────────
\begin{frame}[fragile]{4.12 — AFC avec FactoMineR (et ca)}
\begin{lstlisting}
library(FactoMineR); library(factoextra)

# Tableau de contingence (CSP x Diplome)
tab <- as.table(rbind(
  Cadre   = c(Bac=12,  Lic=65,  Mast=180, Doct=95),
  ProfInt = c(Bac=88,  Lic=110, Mast=50,  Doct=5),
  Employe = c(Bac=230, Lic=45,  Mast=12,  Doct=1),
  Ouvrier = c(Bac=195, Lic=8,   Mast=2,   Doct=0)))

# Test d'independance du chi2
chisq.test(tab)

# Analyse Factorielle des Correspondances
res <- CA(tab, graph = FALSE)

# Valeurs propres et inertie
get_eigenvalue(res)               # lambda_k, % d'inertie
fviz_eig(res, addlabels = TRUE)   # scree plot

# Representation simultanee (biplot)
fviz_ca_biplot(res, repel = TRUE)

# cos2 et contributions
res$row$cos2; res$col$cos2
res$row$contrib; res$col$contrib

# Decomposition du chi2 par cellule
chisq.test(tab)$residuals^2       # ~ contributions au chi2
\end{lstlisting}
\end{frame}

% ────────────────────────────────────────────────────────────
\subsection{4.13 Récapitulatif}
% ────────────────────────────────────────────────────────────
\begin{frame}{4.13 — Récapitulatif du chapitre}
\begin{columns}[T]
\column{0.49\textwidth}
\textbf{À retenir}
\begin{itemize}\small
  \item AFC $=$ ACP des \alert{profils} avec métrique $\chi^2$
  \item Inertie totale $\Phi^2 = \chi^2 / n$
  \item Lignes et colonnes : mêmes valeurs propres
  \item Représentation \alert{simultanée} sur un biplot
  \item Formules de transition $\Rightarrow$ lecture barycentrique
  \item Reconstitution comme \emph{somme} d'écarts à l'indépendance
\end{itemize}

\column{0.49\textwidth}
\begin{block}{Garanties théoriques}
\begin{itemize}\small
  \item SVD généralisée : optimalité (Eckart-Young)
  \item Conservation de l'inertie totale
  \item Décomposition orthogonale du $\chi^2$
\end{itemize}
\end{block}

\begin{alertblock}{Domaine d'application}
\small AFC $\to$ \alert{2 variables qualitatives}.
Pour $p > 2$ variables qualitatives, on utilise l'\textbf{ACM}
(Analyse des Correspondances Multiples), chap. 5.
\end{alertblock}
\end{columns}
\end{frame}
